\section{Pre-Lab Assignment}

Before the in-class session, complete the following:

\begin{enumerate}

  \item \textbf{Review the Lecture Slides.}

  Re-read the slides from the unit on redundancy and diversity in software engineering.
  In particular, review:

  \begin{itemize}
    \item The probabilistic argument for reliability improvement under independence.
    \item The derivation for 1oo2 voting.
    \item IEC~61508 Part~3, Table~A.2, row~3d (Design diversity).
    \item The summary of the Knight \& Leveson experimental evaluation.
  \end{itemize}

  \item \textbf{Derive the 2oo3 Reliability Formula.}

  Suppose each version fails independently with probability $p$ on a given demand.
  Let $q = 1 - p$ denote the probability of success.

  In lecture, we derived the system failure probability for the 1-out-of-2 (1oo2) case, where the system is correct if at least one version is correct:

  \[
    P(\text{system failure}) = p^2.
  \]

  Now consider a 2oo3 (two-out-of-three) voter, \ie the system is correct if at least two of the three versions are correct.

  \begin{itemize}
    \item Enumerate the failure cases for the voter.
    \item Use independence to compute the probability of each case.
    \item Derive an expression for:
    \[
      P(\text{system failure})
    \]
    as a function of $p$ and $q$.
  \end{itemize}

  Your final expression should be written in closed form (no summation required at this stage).
  Show your reasoning clearly.

  \begin{tcolorbox}[colback=yellow!10,colframe=black,boxrule=0.5pt]
  \textbf{Deliverable:} Give your derivation of the 2oo3 system failure probability.
  Your answer must:
  \begin{itemize}
      \item Clearly enumerate the failure cases,
      \item Use the independence assumption explicitly,
      \item Present the final expression as a simplified closed-form formula.
  \end{itemize}

  \medskip
  \textbf{Solution.}
  \begin{itemize}
      \item \textbf{Failure cases:} The 2oo3 system fails when fewer than two versions are correct, \ie when exactly 0 or exactly 1 version is correct.
      \item \textbf{Under independence:} Let $p$ be the probability a version fails and $q = 1 - p$ the probability it is correct. By independence:
      \begin{itemize}
          \item $P(0 \text{ correct}) = P(\text{all three fail}) = p \cdot p \cdot p = p^3$.
          \item $P(1 \text{ correct}) = 3 p^2 q$ (only V1 correct, or only V2, or only V3; each has probability $p^2 q$).
      \end{itemize}
      \item \textbf{Closed form:}
      \[
      P(\text{system failure}) = p^3 + 3p^2 q = p^2(p + 3q) = p^2(3 - 2p).
      \]
  \end{itemize}
  \end{tcolorbox} 
  
  \item \textbf{Set Up Your LLM Environment.}

  Ensure that you have access to an LLM of your choice that you will use to generate program implementations for this lab.

  You may use any of the following:

  \begin{itemize}
    \item \textbf{Gemini Pro}, available to you as a student.
    \item \textbf{Microsoft Copilot}, to which Purdue provides access.
    \item An \textbf{open-weight model} running locally on your laptop (e.g., via a system such as \texttt{ollama}).
  \end{itemize}

  Regardless of the system you choose, confirm that you can:

  \begin{itemize}
    \item Start a fresh session (to create independent ``versions'').
    \item Provide a specification and obtain generated source code.
    \item Run the generated code locally against a simple test suite.
  \end{itemize}

  \item \textbf{Review the Experimental Structure.}

  Review the experimental plan described in the lab handout:

  \begin{itemize}
    \item A frozen specification will be provided.
    \item You will generate $N$ distinct implementations.
    \item A visible test suite will be available during development.
    \item A reserved test set will be used for final evaluation.
  \end{itemize}

  Make sure you understand:

  \begin{itemize}
    \item What is being measured (marginal failure, joint failure, voting reliability).
    \item Where the independence assumption enters the reliability argument.
  \end{itemize}

\end{enumerate}
